% =============================================================================
% 2.2 Problématique et approche générale
% =============================================================================
\section{Problématique et approche générale}
\label{sec:ch2-problematique}

\subsection{Problématique métier}

La question centrale, déjà posée en introduction générale, se reformule ici
en termes opératoires~:

\begin{btpquote}
Comment concevoir et développer une plateforme permettant de capitaliser le
patrimoine documentaire d'un bureau d'études du BTP et de générer
automatiquement des dossiers de réponse aux appels d'offres en s'appuyant
sur ce patrimoine, tout en garantissant la traçabilité des sources, la
conformité réglementaire et la validation humaine des livrables~?
\end{btpquote}

Trois tensions structurent la réponse, et chacune interdit une solution
naïve.

\begin{description}
  \item[Couverture contre précision.] Trop peu d'extraits et l'on rate une
        clause~; trop d'extraits et le modèle se perd dans un contexte
        bruité.
  \item[Texte libre contre structure.] Une partie du savoir est narrative
        (méthodologie d'exécution)~; une autre est tabulaire (moyens,
        montants, certifications). Un traitement uniforme échoue sur l'une
        ou l'autre.
  \item[Autonomie contre responsabilité.] Une offre engage l'entreprise.
        Toute architecture qui délègue la décision à un agent autonome
        déplace la responsabilité sans l'assumer.
\end{description}

\subsection{Approche proposée}

L'approche générale repose sur cinq piliers, qui découlent directement de
l'état de l'art du chapitre~\ref{chap:etat-art}.

\begin{enumerate}
  \item \btpterm{Capitalisation structurée.} Transformer le patrimoine
        documentaire en une base de connaissances interrogeable, en
        préservant la hiérarchie des pièces.
  \item \btpterm{Recherche hybride.} Combiner recherche lexicale et
        recherche sémantique, fusionner par rang réciproque, puis
        ré-ordonner.
  \item \btpterm{Génération contrôlée.} Produire des sections, non des
        documents entiers, sous schéma contraint, avec citations et
        abstentions explicites.
  \item \btpterm{Workflow validé.} Orchestrer le processus avec
        interruption obligatoire pour la revue humaine.
  \item \btpterm{Architecture modulaire.} Concevoir par composants
        indépendants, afin qu'une brique --- le graphe, un n\oe ud
        agentique --- puisse s'ajouter sans reconstruire le reste.
\end{enumerate}

\begin{btpnote}[Ce que le système n'est pas]
Le système visé n'est pas un assistant conversationnel qui rédige un appel
d'offres d'un seul tenant. Ce n'est pas non plus un essaim d'agents
autonomes. C'est une plateforme de capitalisation et de génération
assistée, organisée en quatre capacités~: mémoire documentaire, intelligence
de l'appel d'offres, espace de travail assisté, génération documentaire
versionnée.
\end{btpnote}
